\documentclass[11pt]{article}
\usepackage[margin=1in]{geometry}
\usepackage{amsmath,amssymb,amsthm,mathtools}
\usepackage{graphicx}
\usepackage{microtype}
\usepackage[hidelinks]{hyperref}

\newtheorem{theorem}{Theorem}
\newtheorem{corollary}[theorem]{Corollary}
\newtheorem{lemma}[theorem]{Lemma}
\newcommand{\Rn}{\mathbb{R}^n}
\newcommand{\Linfty}{L^\infty}
\newcommand{\norm}[1]{\left\lVert #1\right\rVert}
\newcommand{\trans}{\tau}

\title{Endpoint products in distributional fractional Sobolev spaces:\\
a strict-loss theorem and a local Hardy--Sobolev subcase}
\author{Candidate partial result for arXiv:2111.13942}
\date{12 August 2026}

\begin{document}
\maketitle

\noindent\textbf{Status.}
Candidate substantial partial result, likely valid; human review is needed.
The exact endpoint algebra question remains open.

\section{Source question}

Comi and Stefani ask in Remark~4.6 of
\emph{Leibniz rules and Gauss--Green formulas in distributional fractional
spaces} \cite{ComiStefani} whether, for $0<\alpha<1$,
\[
 f,g\in S^{\alpha,1}(\Rn)\cap\Linfty(\Rn)
 \quad\Longrightarrow\quad
 \nabla^\alpha_{\mathrm{NL},w}(f,g)\in L^1(\Rn;\mathbb R^n).
\]
The same remark proves the exact equivalence
\[
 \nabla^\alpha_{\mathrm{NL},w}(f,g)\in L^1
 \quad\Longleftrightarrow\quad
 fg\in S^{\alpha,1}(\Rn).                                      \tag{1.1}
\]
Thus the question is precisely whether
$S^{\alpha,1}\cap L^\infty$ is an algebra.

\begin{figure}[ht]
 \centering
 \includegraphics[width=\textwidth]{figures/open_problem_crop.png}
 \caption{Remark~4.6 and the open implication in the source paper
 \cite{ComiStefani}, PDF page 30.}
\end{figure}

Here $S^{\alpha,1}$ consists of $u\in L^1$ whose distributional Riesz
fractional gradient $\nabla^\alpha u$ belongs to $L^1$.  We use
\[
 [u]_{W^{\beta,1}}
 :=\int_{\Rn}\int_{\Rn}
 \frac{|u(x+h)-u(x)|}{|h|^{n+\beta}}\,dx\,dh.             \tag{1.2}
\]

\section{Main partial result}

\begin{theorem}[Every strict lower order is stable]\label{thm:loss}
Let $0<\beta<\alpha<1$ and
$f,g\in S^{\alpha,1}(\Rn)\cap L^\infty(\Rn)$.  Then
\[
 f,g,fg\in W^{\beta,1}(\Rn)\cap L^\infty(\Rn)
       \subset S^{\beta,1}(\Rn)\cap L^\infty(\Rn).
\]
Set
\[
 \mathcal A_{\alpha,\beta}(u)
 :=\frac{\gamma_{n,\alpha}}{\alpha-\beta}
       \norm{\nabla^\alpha u}_{L^1}
   +\frac{2}{\beta}\norm{u}_{L^1},                         \tag{2.1}
\]
where $\gamma_{n,\alpha}$ is the translation constant in
\cite[Proposition~A.1]{BrueCalziComiStefani}.  Then
\begin{align}
 [fg]_{W^{\beta,1}}
 &\le n\omega_n\bigl(
       \norm{g}_{L^\infty}\mathcal A_{\alpha,\beta}(f)
      +\norm{f}_{L^\infty}\mathcal A_{\alpha,\beta}(g)\bigr), \tag{2.2}\\
 \norm{\nabla^\beta(fg)}_{L^1}
 &\le \mu_{n,\beta}n\omega_n\bigl(
       \norm{g}_{L^\infty}\mathcal A_{\alpha,\beta}(f)
      +\norm{f}_{L^\infty}\mathcal A_{\alpha,\beta}(g)\bigr). \tag{2.3}
\end{align}
Consequently $\nabla^\beta_{\mathrm{NL},w}(f,g)\in L^1$ exists and
\[
 \nabla^\beta(fg)=g\nabla^\beta f+f\nabla^\beta g
                  +\nabla^\beta_{\mathrm{NL},w}(f,g)
 \quad\text{in }L^1.                                      \tag{2.4}
\]
\end{theorem}

\begin{theorem}[Exact-order local Hardy--Sobolev subcase]\label{thm:hardy}
Let $0<\alpha<1$ and let
\[
 hS^{\alpha,1}(\Rn):=F^\alpha_{1,2}(\Rn)
\]
be the inhomogeneous Triebel--Lizorkin, or local Hardy--Sobolev, space.
Then
\[
 F^\alpha_{1,2}(\Rn)\hookrightarrow S^{\alpha,1}(\Rn),    \tag{2.5}
\]
and $F^\alpha_{1,2}\cap L^\infty$ is an algebra.  In particular, if
$f,g\in F^\alpha_{1,2}\cap L^\infty$, then the exact-order implication in
(1.1) holds.  More quantitatively,
\[
 \norm{fg}_{F^\alpha_{1,2}}
 \le C\bigl(
    \norm f_{F^\alpha_{1,2}}\norm g_{L^\infty}
   +\norm g_{F^\alpha_{1,2}}\norm f_{L^\infty}\bigr).      \tag{2.6}
\]
This local subclass strictly contains the global Hardy--Sobolev space
$(I-\Delta)^{-\alpha/2}H^1$ used in \cite{BrueCalziComiStefani}.
\end{theorem}

\section{Idea of the proof}

The full hypotheses already give an endpoint translation modulus:
$S^{\alpha,1}$ embeds into $B^\alpha_{1,\infty}$.  Integrating that modulus
against $|h|^{-n-\beta}$ is finite at every $\beta<\alpha$, and the elementary
product difference inequality preserves it.  This proves the full-class
strict-loss theorem without invoking an endpoint Riesz-transform estimate.

At the exact order, the smaller local Hardy--Sobolev scale has precisely the
Littlewood--Paley cancellation needed for pointwise multiplication.  Sickel's
multilinear theorem supplies the algebra estimate.  A low/high-frequency
split then shows that its fractional gradient is integrable, placing the
product inside the source's $S^{\alpha,1}$ space.

The open endpoint is not obtained by taking $\beta\uparrow\alpha$: the bound
has the sharp visible loss $(\alpha-\beta)^{-1}$, and exact
$B^\alpha_{1,\infty}$ regularity alone does not recover the missing
Riesz-transform cancellation.

\section{Proofs}

\begin{proof}[Proof of Theorem~\ref{thm:loss}]
The translation estimate of Bru\`e--Calzi--Comi--Stefani
\cite[Proposition~A.1]{BrueCalziComiStefani}, at $p=1$, is
\[
 \norm{\trans_hu-u}_{L^1}
 \le \gamma_{n,\alpha}|h|^\alpha
      \norm{\nabla^\alpha u}_{L^1}.                       \tag{4.1}
\]
For $|h|>1$ we also have
$\norm{\trans_hu-u}_{L^1}\le2\norm u_{L^1}$.  Polar integration gives
\begin{align*}
 [u]_{W^{\beta,1}}
 &=\int_{\Rn}\frac{\norm{\trans_hu-u}_{L^1}}
                         {|h|^{n+\beta}}\,dh\\
 &\le n\omega_n\left(
    \frac{\gamma_{n,\alpha}}{\alpha-\beta}
          \norm{\nabla^\alpha u}_{L^1}
   +\frac{2}{\beta}\norm u_{L^1}\right)
 =n\omega_n\mathcal A_{\alpha,\beta}(u).                 \tag{4.2}
\end{align*}

Pointwise,
\[
 |(fg)(x+h)-(fg)(x)|
 \le |g(x+h)|\,|f(x+h)-f(x)|
      +|f(x)|\,|g(x+h)-g(x)|.
\]
After integration in $x$ and $h$, (4.2) gives (2.2).  The singular-integral
definition is absolutely convergent for $u\in W^{\beta,1}$ and yields
\[
 \norm{\nabla^\beta u}_{L^1}
 \le\mu_{n,\beta}[u]_{W^{\beta,1}},                        \tag{4.3}
\]
which proves (2.3) and the asserted inclusions.  Applying the same argument
to each input places $f,g$ in $S^{\beta,1}\cap L^\infty$.  Now the equivalence
(1.1), with $\alpha$ replaced by $\beta$, gives (2.4).
\end{proof}

\begin{proof}[Proof of Theorem~\ref{thm:hardy}]
Sickel's Theorem~1 \cite{Sickel}, with
\[
 m=2,\qquad p_1=p_2=p=1,\qquad q=2,\qquad s=\alpha,
\]
applies because
\[
 \alpha>n\max(0,1-1/2-1)=0,
 \qquad \max_j(1/p_j)=1\le1/p=1\le\sum_j1/p_j=2.
\]
It gives (2.6), hence the algebra assertion.

It remains to prove (2.5).  Let $J=(I-\Delta)^{1/2}$.  The lifting property
gives $h:=J^\alpha f\in F^0_{1,2}=h^1$, the local Hardy space.  Since
$h^1\hookrightarrow L^1$ and the Bessel kernel of $J^{-\alpha}$ is in
$L^1$, we obtain $f=J^{-\alpha}h\in L^1$.

For the $j$th fractional-gradient component,
\[
 \nabla^\alpha_j f=T_jh,\qquad
 \widehat{T_jh}(\xi)=m_j(\xi)\widehat h(\xi),
 \quad
 m_j(\xi)=i\xi_j|\xi|^{\alpha-1}(1+|\xi|^2)^{-\alpha/2}.   \tag{4.4}
\]
Choose a smooth cutoff $\chi$ supported near the origin.  The inverse Fourier
transform of $\chi m_j$ is integrable: it is bounded locally and, being a
smooth cutoff of a degree-$\alpha$ homogeneous symbol that vanishes at the
origin, decays as $O(|x|^{-n-\alpha})$.  Thus the low-frequency operator maps
$L^1$ to $L^1$.  The symbol $(1-\chi)m_j$ is a standard order-zero multiplier
supported away from the origin, so the associated local
Calder\'on--Zygmund operator is bounded on $h^1$.  Since $h^1\hookrightarrow
L^1$, both pieces of $T_jh$ lie in $L^1$.  Hence
$\nabla^\alpha f\in L^1$, proving (2.5).

Finally $H^1\hookrightarrow h^1$, so the global space embeds in the local
one.  The inclusion is strict: a nonzero nonnegative $C_c^\infty$ function
belongs to $F^\alpha_{1,2}$, while it cannot belong to
$(I-\Delta)^{-\alpha/2}H^1$ because every global $H^1$ function has integral
zero and $J^\alpha$ preserves the zero-frequency value.
\end{proof}

\section{Verification and limitations}

The proof has three independent consistency checks.
First, the constants in (4.2) are exactly the radial integrals
$\int_0^1r^{\alpha-\beta-1}dr$ and
$\int_1^\infty r^{-\beta-1}dr$.  Second, the product estimate uses only
translations and $L^\infty$ bounds, so no hidden $L^1$ Riesz-transform
estimate enters Theorem~\ref{thm:loss}.  Third, Sickel's parameter conditions
include $p=1$, $q=2$, and every $s=\alpha>0$ exactly as substituted above.

The theorem does \emph{not} establish $fg\in S^{\alpha,1}$ for arbitrary
bounded $S^{\alpha,1}$ inputs.  Nor does it assert that the global
Hardy--Sobolev space is an algebra: if $0\ne f$ is real valued there, then
$\int f^2>0$, whereas every member of that global space has integral zero.
What is proved is that bounded products from the global space land in the
larger local space $F^\alpha_{1,2}$, and hence in $S^{\alpha,1}$.

Eight focused upgrade attempts also tested direct endpoint Kato--Ponce
estimates, the known fractional chain-rule counterexample, accumulating-jump
and lacunary constructions, a closed-graph norm blow-up, and the
$\alpha=0$ obstruction.  None produced a valid full proof or counterexample.
The detailed audit is in the associated attempt note.

Cheap run indexes and bounded searches through 12 August 2026 for the arXiv
id, quoted question, endpoint Leibniz/Kato--Ponce formulations, and the
$S^{\alpha,1}$ algebra problem found no later paper explicitly resolving the
full question.  The local Hardy--Sobolev algebra estimate is known literature
\cite{Sickel}; its implication for the source problem is identified here.
The full-class strict-loss theorem is the promoted new deduction.

\section{Human-review recommendation}

Review (i) the normalization of the double-difference seminorm in (4.2), and
(ii) the low-frequency kernel assertion in (4.4).  If those standard points
are accepted, Theorem~\ref{thm:loss} and the exact-order subclass theorem are
ready to retain as a substantial partial result.  The packet should not be
upgraded to a full solution without a new endpoint cancellation argument or
an explicit counterexample.

\begin{thebibliography}{9}
\bibitem{ComiStefani}
G.~E. Comi and G. Stefani,
\emph{Leibniz rules and Gauss--Green formulas in distributional fractional
spaces}, Rev. Mat. Iberoam. 39 (2023), 1289--1330; arXiv:2111.13942.

\bibitem{BrueCalziComiStefani}
E. Bru\`e, M. Calzi, G.~E. Comi, and G. Stefani,
\emph{A distributional approach to fractional Sobolev spaces and fractional
variation: asymptotics II}, C. R. Math. Acad. Sci. Paris 360 (2022),
589--626; arXiv:2011.03928.

\bibitem{Sickel}
W. Sickel,
\emph{Pointwise multiplication in Triebel--Lizorkin spaces},
Forum Math. 5 (1993), 73--92, Theorem~1.
\end{thebibliography}

\end{document}
